\documentclass{article}
\usepackage{enumitem}
\usepackage{amsmath}
\begin{document}

\title{Chapter 13: Photosynthesis in Higher Plants}
\date{}
\maketitle

\begin{enumerate}[label=Q\arabic*.]

\item The overall equation for oxygenic photosynthesis is:
\\(A) 6CO$_2$ + 6H$_2$O $\xrightarrow{\text{light}}$ C$_6$H$_{12}$O$_6$ + 6O$_2$
\\(B) C$_6$H$_{12}$O$_6$ + 6O$_2$ $\rightarrow$ 6CO$_2$ + 6H$_2$O + ATP
\\(C) 6CO$_2$ + 12H$_2$O $\xrightarrow{\text{light}}$ C$_6$H$_{12}$O$_6$ + 6O$_2$ + 6H$_2$O
\\(D) CO$_2$ + H$_2$O $\rightarrow$ glucose + O$_2$

\item Which photosystem contains P700 as the reaction center chlorophyll?
\\(A) Photosystem I
\\(B) Photosystem II
\\(C) Both PS I and PS II
\\(D) The antenna complex

\item The oxygen released during photosynthesis comes from:
\\(A) CO$_2$ splitting
\\(B) Water photolysis in Photosystem II
\\(C) Glucose oxidation
\\(D) Cyclic photophosphorylation

\item The Calvin cycle (C3 pathway) produces which first stable product after CO$_2$ fixation?
\\(A) Oxaloacetate (OAA)
\\(B) 3-Phosphoglycerate (3-PGA)
\\(C) Phosphoenolpyruvate (PEP)
\\(D) Ribulose-1,5-bisphosphate (RuBP)

\item Assertion: Photorespiration decreases the efficiency of photosynthesis in C$_3$ plants.\\
Reason: In photorespiration, Rubisco adds O$_2$ instead of CO$_2$ to RuBP, producing glycolate, which is oxidized losing CO$_2$ without producing ATP.
\\(A) Both true, Reason is correct explanation
\\(B) Both true, Reason is not correct explanation
\\(C) Assertion true, Reason false
\\(D) Assertion false, Reason true

\item In C$_4$ plants, the first product of CO$_2$ fixation in mesophyll cells is:
\\(A) 3-PGA
\\(B) Oxaloacetate (OAA)
\\(C) Malate
\\(D) Phosphoglyceraldehyde (PGAL)

\item Which enzyme is responsible for the primary CO$_2$ fixation in C$_4$ plants in mesophyll cells?
\\(A) Rubisco
\\(B) PEP carboxylase (PEPC)
\\(C) Pyruvate kinase
\\(D) Malate dehydrogenase

\item Non-cyclic photophosphorylation produces:
\\(A) ATP only
\\(B) NADPH only
\\(C) ATP and NADPH (and releases O$_2$)
\\(D) ATP and NADPH without releasing O$_2$

\item The Emerson enhancement effect demonstrated that:
\\(A) Red light alone is sufficient for maximum photosynthesis
\\(B) Two photosystems work cooperatively; combining far-red and red light gives higher efficiency than either alone
\\(C) Blue light is more efficient than red light
\\(D) Photosynthesis stops above 700 nm wavelength

\item For the synthesis of one molecule of glucose in the Calvin cycle, how many molecules of ATP and NADPH are required?
\\(A) 6 ATP and 6 NADPH
\\(B) 12 ATP and 12 NADPH
\\(C) 18 ATP and 12 NADPH
\\(D) 12 ATP and 18 NADPH

\item The Z-scheme of photosynthesis describes:
\\(A) The cyclic flow of electrons in PS I only
\\(B) The non-cyclic electron flow from water through PS II, the ETC, PS I, to NADP$^+$
\\(C) The carbon fixation reactions in the Calvin cycle
\\(D) The movement of protons across thylakoid membrane

\item Kranz anatomy in C$_4$ plants is characterized by:
\\(A) Large bundle sheath cells with numerous chloroplasts surrounding vascular bundles
\\(B) Small bundle sheath cells with no chloroplasts
\\(C) Stomata only on the lower epidermis
\\(D) Absence of mesophyll cells

\item CAM (Crassulacean Acid Metabolism) plants open stomata at night because:
\\(A) CO$_2$ is more available at night
\\(B) It conserves water; CO$_2$ is fixed into organic acids at night and released for Calvin cycle during day
\\(C) Photosynthesis is more efficient at night
\\(D) Temperature is higher at night

\item Match the photosynthetic pigments with their absorption maxima:\\
1. Chlorophyll a $\rightarrow$ (a) 430 nm and 662 nm\\
2. Chlorophyll b $\rightarrow$ (b) 453 nm and 642 nm\\
3. $\beta$-Carotene $\rightarrow$ (c) 480 nm\\
4. Phycoerythrin $\rightarrow$ (d) 550 nm
\\(A) 1-a, 2-b, 3-c, 4-d
\\(B) 1-b, 2-a, 3-d, 4-c
\\(C) 1-c, 2-d, 3-a, 4-b
\\(D) 1-d, 2-c, 3-b, 4-a

\item The chemiosmotic theory explains ATP synthesis in chloroplasts through:
\\(A) Direct phosphorylation of ADP by electron carriers
\\(B) Proton gradient across thylakoid membrane driving ATP synthase (CF$_0$--CF$_1$)
\\(C) Substrate-level phosphorylation in stroma
\\(D) Reverse electron transport

\end{enumerate}
\end{document}
